% Known-good ctexart preamble for Chinese academic documents. XeLaTeX ONLY.
% Compile recipe (Chinese paths break latexmk):
%   cp <chinese-path>.tex /tmp/latex_build/main.tex
%   cd /tmp/latex_build && xelatex -interaction=nonstopmode main.tex  (x3 passes)
%   cp main.pdf <original-dir>/<new-name>.pdf
\documentclass[UTF8,11pt,a4paper]{ctexart}

\usepackage{amsmath,amssymb,amsthm}
\usepackage{booktabs,longtable,array,multirow}
\usepackage{geometry}
\usepackage{graphicx}
\usepackage{algorithm,algpseudocode}
\usepackage{listings}
\usepackage{xcolor}
\usepackage{hyperref}
\usepackage{enumitem}
\usepackage{caption,subcaption}
\usepackage{tabularx}
\usepackage{cite}

\geometry{left=2.5cm,right=2.5cm,top=2.8cm,bottom=2.8cm}
\hypersetup{colorlinks=true,linkcolor=blue!60!black,citecolor=green!50!black,urlcolor=blue!60!black}

\newtheorem{theorem}{定理}[section]
\newtheorem{proposition}[theorem]{命题}
\newtheorem{definition}[theorem]{定义}
\newtheorem{assumption}[theorem]{假设}
\newtheorem{remark}[theorem]{注}

\lstset{
    basicstyle=\ttfamily\small,
    breaklines=true,
    keywordstyle=\color{blue},
    commentstyle=\color{gray},
    stringstyle=\color{green!50!black},
    numbers=left,
    numberstyle=\tiny\color{gray},
    frame=single,
    language=Python,
    showstringspaces=false
}

\newcommand{\R}{\mathbb{R}}
\newcommand{\E}{\mathbb{E}}
\newcommand{\Var}{\text{Var}}
\newcommand{\SE}{\text{SE}}
\DeclareMathOperator*{\argmin}{arg\,min}   % NOT built-in — required before first use
